\section{技術士のコンピテンシー}

\subsection{令和元年度 問2}
\label{q:r01-reg-apt-02}

\noindent\textbf{原典：} sources/R01/aptitude.pdf、PDF物理3頁。

\textbf{公式正答：} 5（single）。

\paragraph{問題の要旨} 列挙した7項目はすべて平成26年の資質能力に含まれる。

\Needspace{4\baselineskip}\paragraph{確認した条件・記述（独自要約）}\mbox{}\par

\noindent\textbf{提示項目：} ['専門的学識', '問題解決', 'マネジメント', '評価', 'コミュニケーション', 'リーダーシップ', '技術者倫理']\par

\noindent\textbf{指定資料：} 平成26年3月7日技術士分科会資料。\par

\paragraph{問われる核心} 技術士のコンピテンシー

\paragraph{必要知識} 平成26年3月の資質能力には専門的学識、問題解決、マネジメント、評価、コミュニケーション、リーダーシップ、技術者倫理がある。

\paragraph{考え方（査読待ち初稿）} 専門的学識と問題解決は技術判断、マネジメントと評価は業務の実行・改善、コミュニケーションとリーダーシップは関係者との協働、技術者倫理は公益と責任を扱う。指定された平成26年資料にはこの7項目がすべて含まれ、数は7である。

\paragraph{正答を選ぶ理由} 列挙した7項目はすべて平成26年の資質能力に含まれる。 指定された平成26年資料にはこの7項目がすべて含まれ、数は7である。

\Needspace{4\baselineskip}\paragraph{選択肢の検討（査読待ち初稿）}\mbox{}\par

\noindent\textbf{1：不適切な解答候補} 3項目。7項目のうち4項目を根拠なく除外している。\par

\noindent\textbf{2：不適切な解答候補} 4項目。7項目のうち3項目を根拠なく除外している。\par

\noindent\textbf{3：不適切な解答候補} 5項目。7項目のうち2項目を根拠なく除外している。\par

\noindent\textbf{4：不適切な解答候補} 6項目。7項目のうち1項目を根拠なく除外している。\par

\noindent\textbf{5：適切な解答候補} 7項目。提示された全7項目が資料の見出しとして載っている。\par

\paragraph{再利用する考え方} 名称の暗記だけでなく、知識の適用、複合問題への対応、資源配分、成果評価、合意形成、倫理判断の仕事へ対応させる。

\paragraph{共通理論} \hyperref[theory:engineer_act:competency]{第\ref*{theory:engineer_act:competency}節}

\par 参考：\cite{mext_competency2014}。

\par 難易度 1（校準初稿）、優先度未評価。

\section{技術士の義務・責務}

\subsection{平成23年度 問1}
\label{q:h23-reg-apt-01}

\noindent\textbf{原典：} sources/H23/aptitude.pdf、PDF物理2頁。

\textbf{公式正答：} 4（single）。

\paragraph{問題の要旨} アから順に×○○○となる。

\Needspace{4\baselineskip}\paragraph{確認した条件・記述（独自要約）}\mbox{}\par

\noindent\textbf{ア：} \{'原典の主張の要約': '守秘対象を所属組織の取引相手だけに限定する', '独立判定': '×', '理由': '業務上知った秘密は取引相手以外についても対象になる'\}\par

\noindent\textbf{イ：} \{'原典の主張の要約': '社会・環境への影響を評価し公衆に危険があれば良心に従って行動する', '独立判定': '○', '理由': '公益確保のために影響を予測する'\}\par

\noindent\textbf{ウ：} \{'原典の主張の要約': '利益相反情報を雇用者・依頼者へ開示する', '独立判定': '○', '理由': '判断をゆがめる可能性を関係者に説明する'\}\par

\noindent\textbf{エ：} \{'原典の主張の要約': '営業機密から独自に調査した情報の権利は当事者の協議による', '独立判定': '○', '理由': '独自調査だけで一方に帰属すると決め付けない'\}\par

\paragraph{問われる核心} 技術士法の義務・責務と適用条件

\paragraph{必要知識} 守秘の対象と存続、公益優先、客観的判断、情報の権利と利益相反。

\paragraph{考え方（査読待ち初稿）} アは×：業務上知った秘密は取引相手以外についても対象になる。 イは○：公益確保のために影響を予測する。 ウは○：判断をゆがめる可能性を関係者に説明する。 エは○：独自調査だけで一方に帰属すると決め付けない。 判定列は×○○○である。

\paragraph{正答を選ぶ理由} アから順に×○○○となる。  判定列は×○○○である。

\Needspace{4\baselineskip}\paragraph{選択肢の検討（査読待ち初稿）}\mbox{}\par

\noindent\textbf{1：不適切な解答候補} アから順に○×○×。相違する箇所はア（○とするが×）、イ（×とするが○）、エ（×とするが○）である。各記述の理由は独立判定に示した。\par

\noindent\textbf{2：不適切な解答候補} アから順に×○×○。相違する箇所はウ（×とするが○）である。各記述の理由は独立判定に示した。\par

\noindent\textbf{3：不適切な解答候補} アから順に○×○○。相違する箇所はア（○とするが×）、イ（×とするが○）である。各記述の理由は独立判定に示した。\par

\noindent\textbf{4：適切な解答候補} アから順に×○○○。全記述の判断に一致する。\par

\noindent\textbf{5：不適切な解答候補} アから順に○○○×。相違する箇所はア（○とするが×）、エ（×とするが○）である。各記述の理由は独立判定に示した。\par

\paragraph{再利用する考え方} 公益確保と守秘を対立する二択にせず、情報の性質・危険・適切な相手への説明を検討する。独自調査で権利が自動的に移るとは限らない。

\paragraph{共通理論} \hyperref[theory:engineer_act:public]{第\ref*{theory:engineer_act:public}節}

\par 参考：\cite{engineer_act_current,cpd_v3}。

\par 難易度 2（校準初稿）、優先度未評価。

\subsection{平成24年度 問1}
\label{q:h24-reg-apt-01}

\noindent\textbf{原典：} sources/H24/aptitude.pdf、PDF物理1頁。

\textbf{公式正答：} 5（single）。

\paragraph{問題の要旨} 登録部門の力量も継続的に高める必要があり、5が不適切である。

\Needspace{4\baselineskip}\paragraph{確認した条件・記述（独自要約）}\mbox{}\par

\noindent\textbf{条文範囲：} 第44～47条の2。\par

\paragraph{問われる核心} 技術士法の義務・責務と適用条件

\paragraph{必要知識} 第44条の信用失墜禁止、第45条の守秘、第45条の2の公益、第46・47条の名称表示、第47条の2の継続的な資質向上。

\paragraph{考え方（査読待ち初稿）} 1～4はそれぞれ第44条、第45条、第45条の2、第46条に対応する。5は登録部門の知識技能を向上の対象から除外しており、第47条の2に反する。

\paragraph{正答を選ぶ理由} 登録部門の力量も継続的に高める必要があり、5が不適切である。 5は登録部門の知識技能を向上の対象から除外しており、第47条の2に反する。

\Needspace{4\baselineskip}\paragraph{選択肢の検討（査読待ち初稿）}\mbox{}\par

\noindent\textbf{1：不適切な解答候補} 信用失墜・集団の不名誉となる行為を禁止。第44条の内容であり適切である。\par

\noindent\textbf{2：不適切な解答候補} 正当な理由なしに業務上の秘密を漏らし又は盗用しない。資格を失った後も同じ。第45条は資格喪失後の存続も定め、適切である。\par

\noindent\textbf{3：不適切な解答候補} 公共の安全・環境その他公益を害さない努力。第45条の2の公益確保の責務であり適切である。\par

\noindent\textbf{4：不適切な解答候補} 名称表示では登録部門を明示し未登録部門を表示しない。第46条の名称表示時の条件に一致する。\par

\noindent\textbf{5：適切な解答候補} 登録部門の力量は十分なので、それ以外だけの知識技能を向上させる。第47条の2はその業務に関する知識技能の向上を求める。取得した専門部門の学習も続けるので不適切である。\par

\paragraph{再利用する考え方} 主語、名称表示をする場合という条件、退職後も残る義務を確認する。条文の名称と努力義務の性質を区別する。

\paragraph{共通理論} \hyperref[theory:engineer_act:duties]{第\ref*{theory:engineer_act:duties}節}

\par 参考：\cite{engineer_act_current}。

\par 難易度 2（校準初稿）、優先度未評価。

\subsection{平成25年度 問1}
\label{q:h25-reg-apt-01}

\noindent\textbf{原典：} sources/H25/aptitude.pdf、PDF物理1頁。

\textbf{公式正答：} 4（single）。

\par\textbf{独立解答と公式正答に相違あり：} 独立解答 5。独立判定ではアを名称表示時という条件の拡張として誤りと読み、不適切4件・候補5となった。公式は候補4（3件）である。アを名称表示を伴う業務という文脈で適切と読むなら公式に一致するが、公式番号に合わせて条件を補わない。イの情報の範囲と正当な理由の例外も含め、当時の説明・Astra査読で解釈を確認する。

\paragraph{問題の要旨} 提示条文に照らして7つの業務上の記述のうち、不適切なものの個数を判定する。名称表示・情報提供・顧客指示・退職後守秘・技術士補の権限が含まれる。

\Needspace{4\baselineskip}\paragraph{確認した条件・記述（独自要約）}\mbox{}\par

\noindent\textbf{ア：} \{'原典の主張の要約': '業務の実施時には登録技術部門を必ず明示する', '独立判定': '×', '理由': '第46条が要求するのは技術士の名称を表示する場合であり、条件を業務実施全般へ広げている'\}\par

\noindent\textbf{イ：} \{'原典の主張の要約': '業務上知った情報は顧客の許可なく第三者へ渡さない', '独立判定': '○', '理由': '設問の業務上の顧客情報を守る原則として適切。ただし法第45条そのものは秘密と正当な理由を条件とする'\}\par

\noindent\textbf{ウ：} \{'原典の主張の要約': 'どのような顧客の指示でもそのまま実施する', '独立判定': '×', '理由': '違法又は公益を害する指示を無条件に受け入れない'\}\par

\noindent\textbf{エ：} \{'原典の主張の要約': '利益相反を説明し情報を開示する', '独立判定': '○', '理由': '利益相反の開示で判断の公正さを確保する'\}\par

\noindent\textbf{オ：} \{'原典の主張の要約': '退職すれば守秘の制約はない', '独立判定': '×', '理由': '第45条の義務は資格を失った後も存続する'\}\par

\noindent\textbf{カ：} \{'原典の主張の要約': '顧客が認めれば技術士補の名称で主体的に業務を行う', '独立判定': '×', '理由': '第47条は技術士補の名称で技術士を補助する範囲に限定する'\}\par

\noindent\textbf{キ：} \{'原典の主張の要約': '社会環境への影響を評価し公衆に危険があれば良心に従って行動する', '独立判定': '○', '理由': '公益確保を実行する態度である'\}\par

\paragraph{問われる核心} 技術士法の義務・責務と適用条件

\paragraph{必要知識} 第44条の信用失墜禁止、第45条の守秘、第45条の2の公益、第46・47条の名称表示、第47条の2の継続的な資質向上。

\paragraph{考え方（査読待ち初稿）} アは×：第46条が要求するのは技術士の名称を表示する場合であり、条件を業務実施全般へ広げている。 イは○：設問の業務上の顧客情報を守る原則として適切。ただし法第45条そのものは秘密と正当な理由を条件とする。 ウは×：違法又は公益を害する指示を無条件に受け入れない。 エは○：利益相反の開示で判断の公正さを確保する。 オは×：第45条の義務は資格を失った後も存続する。 カは×：第47条は技術士補の名称で技術士を補助する範囲に限定する。 キは○：公益確保を実行する態度である。 よって誤りは4件である。

\paragraph{正答を選ぶ理由} 独立した条文の厳密な読解ではア・ウ・オ・カの4件となり候補5である。公式の候補4と一致しないため、原典のアにおける名称表示の条件を査読する。

\Needspace{4\baselineskip}\paragraph{選択肢の検討（査読待ち初稿）}\mbox{}\par

\noindent\textbf{1：不適切な解答候補} 0件。選択肢は0件だが、誤りと判定したアウオカの4件と合わない。\par

\noindent\textbf{2：不適切な解答候補} 1件。選択肢は1件だが、誤りと判定したアウオカの4件と合わない。\par

\noindent\textbf{3：不適切な解答候補} 2件。選択肢は2件だが、誤りと判定したアウオカの4件と合わない。\par

\noindent\textbf{4：不適切な解答候補} 3件。選択肢は3件だが、誤りと判定したアウオカの4件と合わない。\par

\noindent\textbf{5：適切な解答候補} 4件。個別判定を数えると誤りは4件であり、この数に一致する。\par

\paragraph{再利用する考え方} 主語、名称表示をする場合という条件、退職後も残る義務を確認する。条文の名称と努力義務の性質を区別する。

\paragraph{共通理論} \hyperref[theory:engineer_act:duties]{第\ref*{theory:engineer_act:duties}節}

\par 参考：\cite{engineer_act_current}。

\par 難易度 2（校準初稿）、優先度未評価。

\subsection{平成26年度 問1}
\label{q:h26-reg-apt-01}

\noindent\textbf{原典：} sources/H26/aptitude.pdf、PDF物理1頁。

\textbf{公式正答：} 3（single）。

\paragraph{問題の要旨} 第4章の5空欄へ入る語を選ぶ。 正しい組合せは不名誉／義務／公益／責務／資質である。

\Needspace{4\baselineskip}\paragraph{確認した条件・記述（独自要約）}\mbox{}\par

\noindent\textbf{空欄順：} ['ア', 'イ', 'ウ', 'エ', 'オ']\par

\noindent\textbf{正しい語句の照合：} \{'ア': '不名誉', 'イ': '義務', 'ウ': '公益', 'エ': '責務', 'オ': '資質'\}\par

\noindent\textbf{出題時資料：} \par

\paragraph{問われる核心} 技術士法の義務・責務と適用条件

\paragraph{必要知識} 第44条の信用失墜禁止、第45条の守秘、第45条の2の公益、第46・47条の名称表示、第47条の2の継続的な資質向上。

\paragraph{考え方（査読待ち初稿）} 各空欄の主語と見出しから語句を特定する。 条文又は指定資料に一致する語句は、不名誉、義務、公益、責務、資質である。ほかの組合せは下記の相違がある。

\paragraph{正答を選ぶ理由} 第4章の5空欄へ入る語を選ぶ。 正しい組合せは不名誉／義務／公益／責務／資質である。 ほかの組合せは下記の相違がある。

\Needspace{4\baselineskip}\paragraph{選択肢の検討（査読待ち初稿）}\mbox{}\par

\noindent\textbf{1：不適切な解答候補} 不名誉／義務／安全／責務／能力。ウは安全だけでなく公益、オは資質である。\par

\noindent\textbf{2：不適切な解答候補} 不利益／責務／公益／義務／資質。アは不名誉、イは義務、エは責務である。\par

\noindent\textbf{3：適切な解答候補} 不名誉／義務／公益／責務／資質。第44・45・45条の2・47条の2の見出しと本文に一致する。\par

\noindent\textbf{4：不適切な解答候補} 不名誉／責務／公益／義務／資質。イの義務とエの責務が逆である。\par

\noindent\textbf{5：不適切な解答候補} 不利益／義務／安全／責務／能力。ア・ウ・オが条文と異なる。\par

\paragraph{再利用する考え方} 主語、名称表示をする場合という条件、退職後も残る義務を確認する。条文の名称と努力義務の性質を区別する。

\paragraph{共通理論} \hyperref[theory:engineer_act:duties]{第\ref*{theory:engineer_act:duties}節}

\par 参考：\cite{engineer_act_current}。

\par 難易度 1（校準初稿）、優先度未評価。

\subsection{平成27年度 問1}
\label{q:h27-reg-apt-01}

\noindent\textbf{原典：} sources/H27/aptitude.pdf、PDF物理1頁。

\textbf{公式正答：} 5（single）。

\paragraph{問題の要旨} アから順に×××○○となる。

\Needspace{4\baselineskip}\paragraph{確認した条件・記述（独自要約）}\mbox{}\par

\noindent\textbf{ア：} \{'原典の主張の要約': '依頼者情報から独自調査した情報には守秘が及ばない', '独立判定': '×', '理由': '独自調査という方法だけで秘密性や契約上の権利は消えない'\}\par

\noindent\textbf{イ：} \{'原典の主張の要約': '利害関係者の意見で判断が事実からずれてもやむを得ない', '独立判定': '×', '理由': '他者の意見を聞いても最終判断は客観的事実に基づく'\}\par

\noindent\textbf{ウ：} \{'原典の主張の要約': '登録部門以外を重点的に学べばよい', '独立判定': '×', '理由': '業務に関する専門知識技能も常に更新する'\}\par

\noindent\textbf{エ：} \{'原典の主張の要約': '名称表示時は登録部門を明示し未登録部門は表示しない', '独立判定': '○', '理由': '第46条の規定に一致する'\}\par

\noindent\textbf{オ：} \{'原典の主張の要約': '業務で公共安全・環境その他公益を害さない努力をする', '独立判定': '○', '理由': '第45条の2の規定に一致する'\}\par

\paragraph{問われる核心} 技術士法の義務・責務と適用条件

\paragraph{必要知識} 第44条の信用失墜禁止、第45条の守秘、第45条の2の公益、第46・47条の名称表示、第47条の2の継続的な資質向上。

\paragraph{考え方（査読待ち初稿）} アは×：独自調査という方法だけで秘密性や契約上の権利は消えない。 イは×：他者の意見を聞いても最終判断は客観的事実に基づく。 ウは×：業務に関する専門知識技能も常に更新する。 エは○：第46条の規定に一致する。 オは○：第45条の2の規定に一致する。 判定列は×××○○である。

\paragraph{正答を選ぶ理由} アから順に×××○○となる。  判定列は×××○○である。

\Needspace{4\baselineskip}\paragraph{選択肢の検討（査読待ち初稿）}\mbox{}\par

\noindent\textbf{1：不適切な解答候補} アから順に○○×○○。相違する箇所はア（○とするが×）、イ（○とするが×）である。各記述の理由は独立判定に示した。\par

\noindent\textbf{2：不適切な解答候補} アから順に○×○××。相違する箇所はア（○とするが×）、ウ（○とするが×）、エ（×とするが○）、オ（×とするが○）である。各記述の理由は独立判定に示した。\par

\noindent\textbf{3：不適切な解答候補} アから順に○○○×○。相違する箇所はア（○とするが×）、イ（○とするが×）、ウ（○とするが×）、エ（×とするが○）である。各記述の理由は独立判定に示した。\par

\noindent\textbf{4：不適切な解答候補} アから順に×○×○×。相違する箇所はイ（○とするが×）、オ（×とするが○）である。各記述の理由は独立判定に示した。\par

\noindent\textbf{5：適切な解答候補} アから順に×××○○。全記述の判断に一致する。\par

\paragraph{再利用する考え方} 主語、名称表示をする場合という条件、退職後も残る義務を確認する。条文の名称と努力義務の性質を区別する。

\paragraph{共通理論} \hyperref[theory:engineer_act:duties]{第\ref*{theory:engineer_act:duties}節}

\par 参考：\cite{engineer_act_current}。

\par 難易度 2（校準初稿）、優先度未評価。

\subsection{平成28年度 問1}
\label{q:h28-reg-apt-01}

\noindent\textbf{原典：} sources/H28/aptitude.pdf、PDF物理1頁。

\textbf{公式正答：} 2（single）。

\paragraph{問題の要旨} 正しいと判定する記述はエで、1件である。

\Needspace{4\baselineskip}\paragraph{確認した条件・記述（独自要約）}\mbox{}\par

\noindent\textbf{ア：} \{'原典の主張の要約': 'あらゆる顧客指示をそのまま実施する', '独立判定': '×', '理由': '公益や法令に反する指示を無条件に履行しない'\}\par

\noindent\textbf{イ：} \{'原典の主張の要約': '利害関係者への配慮で事実と異なる判断でもよい', '独立判定': '×', '理由': '判断は客観的な事実に基づく'\}\par

\noindent\textbf{ウ：} \{'原典の主張の要約': '登録部門以外の知識技能を重点的に高める', '独立判定': '×', '理由': '登録済みの領域も資質向上の対象である'\}\par

\noindent\textbf{エ：} \{'原典の主張の要約': '名称表示している専門領域の能力開発により業務領域を広げる', '独立判定': '○', '理由': '技術の進歩に応じた能力開発は業務能力を高める。ただし新部門の名称表示は別途登録を要する'\}\par

\noindent\textbf{オ：} \{'原典の主張の要約': '顧客利益を害する場合は公益より守秘を優先', '独立判定': '×', '理由': '公衆に及ぶ危険を顧客利益だけで隠さない'\}\par

\noindent\textbf{カ：} \{'原典の主張の要約': '顧客が認めれば技術士補の名称で主体的に業務を行う', '独立判定': '×', '理由': '第47条の補助の条件に反する'\}\par

\noindent\textbf{キ：} \{'原典の主張の要約': '退職後の守秘制約はない', '独立判定': '×', '理由': '第45条は資格喪失後も守秘を求める'\}\par

\paragraph{問われる核心} 技術士法の義務・責務と適用条件

\paragraph{必要知識} 第44条の信用失墜禁止、第45条の守秘、第45条の2の公益、第46・47条の名称表示、第47条の2の継続的な資質向上。

\paragraph{考え方（査読待ち初稿）} アは×：公益や法令に反する指示を無条件に履行しない。 イは×：判断は客観的な事実に基づく。 ウは×：登録済みの領域も資質向上の対象である。 エは○：技術の進歩に応じた能力開発は業務能力を高める。ただし新部門の名称表示は別途登録を要する。 オは×：公衆に及ぶ危険を顧客利益だけで隠さない。 カは×：第47条の補助の条件に反する。 キは×：第45条は資格喪失後も守秘を求める。 よって正しいは1件である。

\paragraph{正答を選ぶ理由} 正しいと判定する記述はエで、1件である。  よって正しいは1件である。

\Needspace{4\baselineskip}\paragraph{選択肢の検討（査読待ち初稿）}\mbox{}\par

\noindent\textbf{1：不適切な解答候補} 0件。選択肢は0件だが、正しいと判定したエの1件と合わない。\par

\noindent\textbf{2：適切な解答候補} 1件。個別判定を数えると正しいは1件であり、この数に一致する。\par

\noindent\textbf{3：不適切な解答候補} 2件。選択肢は2件だが、正しいと判定したエの1件と合わない。\par

\noindent\textbf{4：不適切な解答候補} 3件。選択肢は3件だが、正しいと判定したエの1件と合わない。\par

\noindent\textbf{5：不適切な解答候補} 4件。選択肢は4件だが、正しいと判定したエの1件と合わない。\par

\paragraph{再利用する考え方} 主語、名称表示をする場合という条件、退職後も残る義務を確認する。条文の名称と努力義務の性質を区別する。

\paragraph{共通理論} \hyperref[theory:engineer_act:duties]{第\ref*{theory:engineer_act:duties}節}

\par 参考：\cite{engineer_act_current}。

\par 難易度 2（校準初稿）、優先度未評価。

\subsection{平成29年度 問1}
\label{q:h29-reg-apt-01}

\noindent\textbf{原典：} sources/H29/aptitude.pdf、PDF物理1頁。

\textbf{公式正答：} 3（single）。

\paragraph{問題の要旨} 第4章の7空欄へ入る語を選ぶ。 正しい組合せは不名誉／義務／盗用／公益／責務／技術部門／資質である。

\Needspace{4\baselineskip}\paragraph{確認した条件・記述（独自要約）}\mbox{}\par

\noindent\textbf{空欄順：} ['ア', 'イ', 'ウ', 'エ', 'オ', 'カ', 'キ']\par

\noindent\textbf{正しい語句の照合：} \{'ア': '不名誉', 'イ': '義務', 'ウ': '盗用', 'エ': '公益', 'オ': '責務', 'カ': '技術部門', 'キ': '資質'\}\par

\noindent\textbf{出題時資料：} \par

\paragraph{問われる核心} 技術士法の義務・責務と適用条件

\paragraph{必要知識} 第44条の信用失墜禁止、第45条の守秘、第45条の2の公益、第46・47条の名称表示、第47条の2の継続的な資質向上。

\paragraph{考え方（査読待ち初稿）} 各空欄の主語と見出しから語句を特定する。 条文又は指定資料に一致する語句は、不名誉、義務、盗用、公益、責務、技術部門、資質である。ほかの組合せは下記の相違がある。

\paragraph{正答を選ぶ理由} 第4章の7空欄へ入る語を選ぶ。 正しい組合せは不名誉／義務／盗用／公益／責務／技術部門／資質である。 ほかの組合せは下記の相違がある。

\Needspace{4\baselineskip}\paragraph{選択肢の検討（査読待ち初稿）}\mbox{}\par

\noindent\textbf{1：不適切な解答候補} 不名誉／義務／盗用／安全／責務／技術部門／能力。エは公益、キは資質である。\par

\noindent\textbf{2：不適切な解答候補} 信用失墜／責務／盗作／公益／義務／技術部門／資質。アは不名誉、イは義務、ウは盗用、オは責務である。\par

\noindent\textbf{3：適切な解答候補} 不名誉／義務／盗用／公益／責務／技術部門／資質。見出し、守秘の禁止行為、名称表示の登録区分に一致する。\par

\noindent\textbf{4：不適切な解答候補} 不名誉／責務／盗作／公益／義務／専門部門／資質。イ・ウ・オ・カが条文の語句と異なる。\par

\noindent\textbf{5：不適切な解答候補} 信用失墜／義務／盗作／安全／責務／専門部門／能力。ア・ウ・エ・カ・キが条文と異なる。\par

\paragraph{再利用する考え方} 主語、名称表示をする場合という条件、退職後も残る義務を確認する。条文の名称と努力義務の性質を区別する。

\paragraph{共通理論} \hyperref[theory:engineer_act:duties]{第\ref*{theory:engineer_act:duties}節}

\par 参考：\cite{engineer_act_current}。

\par 難易度 1（校準初稿）、優先度未評価。

\subsection{平成29年度 問2}
\label{q:h29-reg-apt-02}

\noindent\textbf{原典：} sources/H29/aptitude.pdf、PDF物理2頁。

\textbf{公式正答：} 5（single）。

\paragraph{問題の要旨} 誤りと判定する記述はイエカキで、4件である。

\Needspace{4\baselineskip}\paragraph{確認した条件・記述（独自要約）}\mbox{}\par

\noindent\textbf{ア：} \{'原典の主張の要約': '社会環境影響を評価し公衆への危険に良心に従って行動する', '独立判定': '○', '理由': '公益のために影響を予測して行動する'\}\par

\noindent\textbf{イ：} \{'原典の主張の要約': '依頼者情報から独自調査した情報は守秘対象外', '独立判定': '×', '理由': '独自調査だけでは秘密性や情報の権利はなくならない'\}\par

\noindent\textbf{ウ：} \{'原典の主張の要約': '部下の企画を承認する際に自身でも安全性・信頼度を検討する', '独立判定': '○', '理由': '承認者も専門家として検討・説明の責任を負う'\}\par

\noindent\textbf{エ：} \{'原典の主張の要約': '顧客判断に安全懸念があっても指摘しない', '独立判定': '×', '理由': '予想される危険を顧客に説明する必要がある'\}\par

\noindent\textbf{オ：} \{'原典の主張の要約': '利益相反情報を説明する', '独立判定': '○', '理由': '判断の公正さを保つため関係を明らかにする'\}\par

\noindent\textbf{カ：} \{'原典の主張の要約': '専門分野外でも自己判断だけで進める', '独立判定': '×', '理由': '自己の力量の範囲を踏まえ専門家の協力を求める'\}\par

\noindent\textbf{キ：} \{'原典の主張の要約': '顧客が認めれば技術士補が技術士に代わり主体的に業務を行う', '独立判定': '×', '理由': '技術士補の補助という役割を越える'\}\par

\paragraph{問われる核心} 技術士法の義務・責務と適用条件

\paragraph{必要知識} 第44条の信用失墜禁止、第45条の守秘、第45条の2の公益、第46・47条の名称表示、第47条の2の継続的な資質向上。

\paragraph{考え方（査読待ち初稿）} アは○：公益のために影響を予測して行動する。 イは×：独自調査だけでは秘密性や情報の権利はなくならない。 ウは○：承認者も専門家として検討・説明の責任を負う。 エは×：予想される危険を顧客に説明する必要がある。 オは○：判断の公正さを保つため関係を明らかにする。 カは×：自己の力量の範囲を踏まえ専門家の協力を求める。 キは×：技術士補の補助という役割を越える。 よって誤りは4件である。

\paragraph{正答を選ぶ理由} 誤りと判定する記述はイエカキで、4件である。  よって誤りは4件である。

\Needspace{4\baselineskip}\paragraph{選択肢の検討（査読待ち初稿）}\mbox{}\par

\noindent\textbf{1：不適切な解答候補} 0件。選択肢は0件だが、誤りと判定したイエカキの4件と合わない。\par

\noindent\textbf{2：不適切な解答候補} 1件。選択肢は1件だが、誤りと判定したイエカキの4件と合わない。\par

\noindent\textbf{3：不適切な解答候補} 2件。選択肢は2件だが、誤りと判定したイエカキの4件と合わない。\par

\noindent\textbf{4：不適切な解答候補} 3件。選択肢は3件だが、誤りと判定したイエカキの4件と合わない。\par

\noindent\textbf{5：適切な解答候補} 4件。個別判定を数えると誤りは4件であり、この数に一致する。\par

\paragraph{再利用する考え方} 主語、名称表示をする場合という条件、退職後も残る義務を確認する。条文の名称と努力義務の性質を区別する。

\paragraph{共通理論} \hyperref[theory:engineer_act:duties]{第\ref*{theory:engineer_act:duties}節}

\par 参考：\cite{engineer_act_current}。

\par 難易度 2（校準初稿）、優先度未評価。

\subsection{平成30年度 問1}
\label{q:h30-reg-apt-01}

\noindent\textbf{原典：} sources/H30/aptitude.pdf、PDF物理1頁。

\textbf{公式正答：} 2（single）。

\paragraph{問題の要旨} 第4章の5空欄へ入る語を選ぶ。 正しい組合せは禁止／義務／公益／責務／制限である。

\Needspace{4\baselineskip}\paragraph{確認した条件・記述（独自要約）}\mbox{}\par

\noindent\textbf{空欄順：} ['ア', 'イ', 'ウ', 'エ', 'オ']\par

\noindent\textbf{正しい語句の照合：} \{'ア': '禁止', 'イ': '義務', 'ウ': '公益', 'エ': '責務', 'オ': '制限'\}\par

\noindent\textbf{出題時資料：} \par

\paragraph{問われる核心} 技術士法の義務・責務と適用条件

\paragraph{必要知識} 第44条の信用失墜禁止、第45条の守秘、第45条の2の公益、第46・47条の名称表示、第47条の2の継続的な資質向上。

\paragraph{考え方（査読待ち初稿）} 各空欄の主語と見出しから語句を特定する。 条文又は指定資料に一致する語句は、禁止、義務、公益、責務、制限である。ほかの組合せは下記の相違がある。

\paragraph{正答を選ぶ理由} 第4章の5空欄へ入る語を選ぶ。 正しい組合せは禁止／義務／公益／責務／制限である。 ほかの組合せは下記の相違がある。

\Needspace{4\baselineskip}\paragraph{選択肢の検討（査読待ち初稿）}\mbox{}\par

\noindent\textbf{1：不適切な解答候補} 制限／責務／利益／義務／制約。5語すべてが見出し・条文と異なる。\par

\noindent\textbf{2：適切な解答候補} 禁止／義務／公益／責務／制限。信用失墜の禁止、秘密保持の義務、公益確保の責務、業務の制限に一致する。\par

\noindent\textbf{3：不適切な解答候補} 禁止／義務／利益／責務／制約。ウは公益、オは制限である。\par

\noindent\textbf{4：不適切な解答候補} 禁止／責務／利益／義務／制限。イ・ウ・エが異なる。\par

\noindent\textbf{5：不適切な解答候補} 制限／責務／公益／義務／制約。ア・イ・エ・オが異なる。\par

\paragraph{再利用する考え方} 主語、名称表示をする場合という条件、退職後も残る義務を確認する。条文の名称と努力義務の性質を区別する。

\paragraph{共通理論} \hyperref[theory:engineer_act:duties]{第\ref*{theory:engineer_act:duties}節}

\par 参考：\cite{engineer_act_current}。

\par 難易度 1（校準初稿）、優先度未評価。

\subsection{平成30年度 問2}
\label{q:h30-reg-apt-02}

\noindent\textbf{原典：} sources/H30/aptitude.pdf、PDF物理2頁。

\textbf{公式正答：} 5（single）。

\paragraph{問題の要旨} 誤りと判定する記述はアイウエオで、5件である。

\Needspace{4\baselineskip}\paragraph{確認した条件・記述（独自要約）}\mbox{}\par

\noindent\textbf{ア：} \{'原典の主張の要約': '営業機密から独自調査した情報の財産権は発注者に属する限りではない', '独立判定': '×', '理由': '独自調査だけで権利関係が当然に切り替わるとはいえず協議が必要'\}\par

\noindent\textbf{イ：} \{'原典の主張の要約': '顧客と公衆の利益が相反すれば所属企業を最優先', '独立判定': '×', '理由': '公衆の安全・健康・福利を優先する'\}\par

\noindent\textbf{ウ：} \{'原典の主張の要約': '退職後の守秘制約はない', '独立判定': '×', '理由': '守秘は退職・資格喪失後も存続する'\}\par

\noindent\textbf{エ：} \{'原典の主張の要約': '顧客が認めれば技術士補の名称で主体的に業務を行う', '独立判定': '×', '理由': '技術士補の名称で行えるのは技術士を補助する場合である'\}\par

\noindent\textbf{オ：} \{'原典の主張の要約': '登録部門以外の知識技能を重点的に高める', '独立判定': '×', '理由': '自分の業務に関する知識技能も高める責務がある'\}\par

\paragraph{問われる核心} 技術士法の義務・責務と適用条件

\paragraph{必要知識} 第44条の信用失墜禁止、第45条の守秘、第45条の2の公益、第46・47条の名称表示、第47条の2の継続的な資質向上。

\paragraph{考え方（査読待ち初稿）} アは×：独自調査だけで権利関係が当然に切り替わるとはいえず協議が必要。 イは×：公衆の安全・健康・福利を優先する。 ウは×：守秘は退職・資格喪失後も存続する。 エは×：技術士補の名称で行えるのは技術士を補助する場合である。 オは×：自分の業務に関する知識技能も高める責務がある。 よって誤りは5件である。

\paragraph{正答を選ぶ理由} 誤りと判定する記述はアイウエオで、5件である。  よって誤りは5件である。

\Needspace{4\baselineskip}\paragraph{選択肢の検討（査読待ち初稿）}\mbox{}\par

\noindent\textbf{1：不適切な解答候補} 1件。選択肢は1件だが、誤りと判定したアイウエオの5件と合わない。\par

\noindent\textbf{2：不適切な解答候補} 2件。選択肢は2件だが、誤りと判定したアイウエオの5件と合わない。\par

\noindent\textbf{3：不適切な解答候補} 3件。選択肢は3件だが、誤りと判定したアイウエオの5件と合わない。\par

\noindent\textbf{4：不適切な解答候補} 4件。選択肢は4件だが、誤りと判定したアイウエオの5件と合わない。\par

\noindent\textbf{5：適切な解答候補} 5件。個別判定を数えると誤りは5件であり、この数に一致する。\par

\paragraph{再利用する考え方} 主語、名称表示をする場合という条件、退職後も残る義務を確認する。条文の名称と努力義務の性質を区別する。

\paragraph{共通理論} \hyperref[theory:engineer_act:duties]{第\ref*{theory:engineer_act:duties}節}

\par 参考：\cite{engineer_act_current}。

\par 難易度 2（校準初稿）、優先度未評価。

\subsection{令和元年度 問1}
\label{q:r01-reg-apt-01}

\noindent\textbf{原典：} sources/R01/aptitude.pdf、PDF物理1頁。

\textbf{公式正答：} 5（single）。

\paragraph{問題の要旨} 第4章の7空欄へ入る語を選ぶ。 正しい組合せは義務／公益／責務／技術部門／制限／準用／資質である。

\Needspace{4\baselineskip}\paragraph{確認した条件・記述（独自要約）}\mbox{}\par

\noindent\textbf{空欄順：} ['ア', 'イ', 'ウ', 'エ', 'オ', 'カ', 'キ']\par

\noindent\textbf{正しい語句の照合：} \{'ア': '義務', 'イ': '公益', 'ウ': '責務', 'エ': '技術部門', 'オ': '制限', 'カ': '準用', 'キ': '資質'\}\par

\noindent\textbf{出題時資料：} \par

\paragraph{問われる核心} 技術士法の義務・責務と適用条件

\paragraph{必要知識} 第44条の信用失墜禁止、第45条の守秘、第45条の2の公益、第46・47条の名称表示、第47条の2の継続的な資質向上。

\paragraph{考え方（査読待ち初稿）} 各空欄の主語と見出しから語句を特定する。 条文又は指定資料に一致する語句は、義務、公益、責務、技術部門、制限、準用、資質である。ほかの組合せは下記の相違がある。

\paragraph{正答を選ぶ理由} 第4章の7空欄へ入る語を選ぶ。 正しい組合せは義務／公益／責務／技術部門／制限／準用／資質である。 ほかの組合せは下記の相違がある。

\Needspace{4\baselineskip}\paragraph{選択肢の検討（査読待ち初稿）}\mbox{}\par

\noindent\textbf{1：不適切な解答候補} 義務／公益／責務／技術部門／制限／準用／能力。キは能力ではなく資質である。\par

\noindent\textbf{2：不適切な解答候補} 責務／安全／義務／専門部門／制約／適用／能力。ア・イ・ウ・エ・オ・カ・キが条文と異なる。\par

\noindent\textbf{3：不適切な解答候補} 義務／公益／責務／技術部門／制約／適用／資質。オは制限、カは準用である。\par

\noindent\textbf{4：不適切な解答候補} 責務／安全／義務／専門部門／制約／準用／資質。ア・イ・ウ・エ・オが条文と異なる。\par

\noindent\textbf{5：適切な解答候補} 義務／公益／責務／技術部門／制限／準用／資質。秘密保持・名称表示は義務、公益・資質向上は責務。登録技術部門、業務の制限、準用という語も一致する。\par

\paragraph{再利用する考え方} 主語、名称表示をする場合という条件、退職後も残る義務を確認する。条文の名称と努力義務の性質を区別する。

\paragraph{共通理論} \hyperref[theory:engineer_act:duties]{第\ref*{theory:engineer_act:duties}節}

\par 参考：\cite{engineer_act_current}。

\par 難易度 1（校準初稿）、優先度未評価。

\subsection{令和元年度再試験 問1}
\label{q:r01-ret-apt-01}

\noindent\textbf{原典：} sources/R01\_retest/aptitude.pdf、PDF物理1頁。

\textbf{公式正答：} 3（single）。

\paragraph{問題の要旨} 適性科目の目的と第4章の規範としての性格を選ぶ。 正しい組合せは社会規範／法規範／倫理規範である。

\Needspace{4\baselineskip}\paragraph{確認した条件・記述（独自要約）}\mbox{}\par

\noindent\textbf{空欄順：} ['ア', 'イ', 'ウ']\par

\noindent\textbf{正しい語句の照合：} \{'ア': '社会規範', 'イ': '法規範', 'ウ': '倫理規範'\}\par

\noindent\textbf{出題時資料：} \par

\paragraph{問われる核心} 技術士法の義務・責務と適用条件

\paragraph{必要知識} 社会規範の下位概念としての法規範と倫理規範。

\paragraph{考え方（査読待ち初稿）} 各空欄の主語と見出しから語句を特定する。 条文又は指定資料に一致する語句は、社会規範、法規範、倫理規範である。ほかの組合せは下記の相違がある。

\paragraph{正答を選ぶ理由} 適性科目の目的と第4章の規範としての性格を選ぶ。 正しい組合せは社会規範／法規範／倫理規範である。 ほかの組合せは下記の相違がある。

\Needspace{4\baselineskip}\paragraph{選択肢の検討（査読待ち初稿）}\mbox{}\par

\noindent\textbf{1：不適切な解答候補} 社会規範／倫理規範／法規範。イとウが逆である。法として制定された第4章は形式上法規範である。\par

\noindent\textbf{2：不適切な解答候補} 行動規範／法規範／倫理規範。アは法と倫理を包括する社会規範である。\par

\noindent\textbf{3：適切な解答候補} 社会規範／法規範／倫理規範。法・倫理を包括する社会規範、形式上の法、内容上の倫理の関係に一致する。\par

\noindent\textbf{4：不適切な解答候補} 行動規範／倫理規範／行動規範。ア・イ・ウの区別が崩れている。\par

\noindent\textbf{5：不適切な解答候補} 社会規範／行動規範／倫理規範。イは一般的行動規範より具体的な法規範である。\par

\paragraph{再利用する考え方} 形式として法であることと、内容として職業倫理を表すことは両立する。

\paragraph{共通理論} \hyperref[theory:engineer_act:norms]{第\ref*{theory:engineer_act:norms}節}

\par 参考：\cite{engineer_act_current}。

\par 難易度 1（校準初稿）、優先度未評価。

\subsection{令和元年度再試験 問2}
\label{q:r01-ret-apt-02}

\noindent\textbf{原典：} sources/R01\_retest/aptitude.pdf、PDF物理2頁。

\textbf{公式正答：} 3（single）。

\paragraph{問題の要旨} 正しいと判定する記述はウカで、2件である。

\Needspace{4\baselineskip}\paragraph{確認した条件・記述（独自要約）}\mbox{}\par

\noindent\textbf{ア：} \{'原典の主張の要約': '登録部門以外の知識技能を重点的に向上', '独立判定': '×', '理由': '業務に関する登録領域も継続して学ぶ'\}\par

\noindent\textbf{イ：} \{'原典の主張の要約': 'あらゆる顧客指示を守秘優先で履行', '独立判定': '×', '理由': '公益や法令に反する命令まで無条件に履行しない'\}\par

\noindent\textbf{ウ：} \{'原典の主張の要約': '業務に関する知識技能と専門領域の能力を向上', '独立判定': '○', '理由': '第47条の2の資質向上に沿う'\}\par

\noindent\textbf{エ：} \{'原典の主張の要約': '利害関係者への配慮で事実から異なる判断も容認', '独立判定': '×', '理由': '客観的事実に基づいて判断する'\}\par

\noindent\textbf{オ：} \{'原典の主張の要約': '顧客利益を害する場合は公益より守秘を優先', '独立判定': '×', '理由': '顧客の利益だけで公衆の危険を放置しない'\}\par

\noindent\textbf{カ：} \{'原典の主張の要約': '資格を失った後も守秘義務を守る', '独立判定': '○', '理由': '第45条後段の規定に沿う'\}\par

\noindent\textbf{キ：} \{'原典の主張の要約': '顧客が認めれば技術士補の名称で主体的に業務を行う', '独立判定': '×', '理由': '第47条は補助としての名称使用を求める'\}\par

\paragraph{問われる核心} 技術士法の義務・責務と適用条件

\paragraph{必要知識} 第44条の信用失墜禁止、第45条の守秘、第45条の2の公益、第46・47条の名称表示、第47条の2の継続的な資質向上。

\paragraph{考え方（査読待ち初稿）} アは×：業務に関する登録領域も継続して学ぶ。 イは×：公益や法令に反する命令まで無条件に履行しない。 ウは○：第47条の2の資質向上に沿う。 エは×：客観的事実に基づいて判断する。 オは×：顧客の利益だけで公衆の危険を放置しない。 カは○：第45条後段の規定に沿う。 キは×：第47条は補助としての名称使用を求める。 よって正しいは2件である。

\paragraph{正答を選ぶ理由} 正しいと判定する記述はウカで、2件である。  よって正しいは2件である。

\Needspace{4\baselineskip}\paragraph{選択肢の検討（査読待ち初稿）}\mbox{}\par

\noindent\textbf{1：不適切な解答候補} 0件。選択肢は0件だが、正しいと判定したウカの2件と合わない。\par

\noindent\textbf{2：不適切な解答候補} 1件。選択肢は1件だが、正しいと判定したウカの2件と合わない。\par

\noindent\textbf{3：適切な解答候補} 2件。個別判定を数えると正しいは2件であり、この数に一致する。\par

\noindent\textbf{4：不適切な解答候補} 3件。選択肢は3件だが、正しいと判定したウカの2件と合わない。\par

\noindent\textbf{5：不適切な解答候補} 4件。選択肢は4件だが、正しいと判定したウカの2件と合わない。\par

\paragraph{再利用する考え方} 主語、名称表示をする場合という条件、退職後も残る義務を確認する。条文の名称と努力義務の性質を区別する。

\paragraph{共通理論} \hyperref[theory:engineer_act:duties]{第\ref*{theory:engineer_act:duties}節}

\par 参考：\cite{engineer_act_current}。

\par 難易度 2（校準初稿）、優先度未評価。

\subsection{令和2年度 問1}
\label{q:r02-reg-apt-01}

\noindent\textbf{原典：} sources/R02/aptitude.pdf、PDF物理1頁。

\textbf{公式正答：} 5（single）。

\paragraph{問題の要旨} 第4章の7空欄へ入る語を選ぶ。 正しい組合せは義務／公益／責務／技術部門／制限／準用／資質である。

\Needspace{4\baselineskip}\paragraph{確認した条件・記述（独自要約）}\mbox{}\par

\noindent\textbf{空欄順：} ['ア', 'イ', 'ウ', 'エ', 'オ', 'カ', 'キ']\par

\noindent\textbf{正しい語句の照合：} \{'ア': '義務', 'イ': '公益', 'ウ': '責務', 'エ': '技術部門', 'オ': '制限', 'カ': '準用', 'キ': '資質'\}\par

\noindent\textbf{出題時資料：} \par

\paragraph{問われる核心} 技術士法の義務・責務と適用条件

\paragraph{必要知識} 第44条の信用失墜禁止、第45条の守秘、第45条の2の公益、第46・47条の名称表示、第47条の2の継続的な資質向上。

\paragraph{考え方（査読待ち初稿）} 各空欄の主語と見出しから語句を特定する。 条文又は指定資料に一致する語句は、義務、公益、責務、技術部門、制限、準用、資質である。ほかの組合せは下記の相違がある。

\paragraph{正答を選ぶ理由} 第4章の7空欄へ入る語を選ぶ。 正しい組合せは義務／公益／責務／技術部門／制限／準用／資質である。 ほかの組合せは下記の相違がある。

\Needspace{4\baselineskip}\paragraph{選択肢の検討（査読待ち初稿）}\mbox{}\par

\noindent\textbf{1：不適切な解答候補} 義務／公益／責務／技術部門／制限／準用／能力。キは能力ではなく資質である。\par

\noindent\textbf{2：不適切な解答候補} 責務／安全／義務／専門部門／制約／適用／能力。ア・イ・ウ・エ・オ・カ・キが条文と異なる。\par

\noindent\textbf{3：不適切な解答候補} 義務／公益／責務／技術部門／制約／適用／資質。オは制限、カは準用である。\par

\noindent\textbf{4：不適切な解答候補} 責務／安全／義務／専門部門／制約／準用／資質。ア・イ・ウ・エ・オが条文と異なる。\par

\noindent\textbf{5：適切な解答候補} 義務／公益／責務／技術部門／制限／準用／資質。秘密保持・名称表示は義務、公益・資質向上は責務。登録技術部門、業務の制限、準用という語も一致する。\par

\paragraph{再利用する考え方} 主語、名称表示をする場合という条件、退職後も残る義務を確認する。条文の名称と努力義務の性質を区別する。

\paragraph{共通理論} \hyperref[theory:engineer_act:duties]{第\ref*{theory:engineer_act:duties}節}

\par 参考：\cite{engineer_act_current}。

\par 難易度 1（校準初稿）、優先度未評価。

\subsection{令和3年度 問1}
\label{q:r03-reg-apt-01}

\noindent\textbf{原典：} sources/R03/aptitude.pdf、PDF物理1頁。

\textbf{公式正答：} 3（single）。

\paragraph{問題の要旨} 誤りと判定する記述はイウエオカで、5件である。

\Needspace{4\baselineskip}\paragraph{確認した条件・記述（独自要約）}\mbox{}\par

\noindent\textbf{ア：} \{'原典の主張の要約': '業務上知った情報を顧客の許可なく第三者に提供しない', '独立判定': '○', '理由': '顧客の業務情報を保護する原則として適切。第45条の正当な理由という例外まで否定するものとしては読まない'\}\par

\noindent\textbf{イ：} \{'原典の主張の要約': '退職後の守秘制約はない', '独立判定': '×', '理由': '守秘は退職・資格喪失後も存続する'\}\par

\noindent\textbf{ウ：} \{'原典の主張の要約': 'どのような顧客指示でもそのまま履行する', '独立判定': '×', '理由': '公益と法令に反する指示まで受け入れない'\}\par

\noindent\textbf{エ：} \{'原典の主張の要約': '顧客利益を害する場合は公益より守秘を優先', '独立判定': '×', '理由': '公衆の危険を顧客利益だけで放置しない'\}\par

\noindent\textbf{オ：} \{'原典の主張の要約': '名称表示で登録部門を明示するが技術士補には同じ制限がない', '独立判定': '×', '理由': '第47条2項は名称表示の第46条を技術士補へ準用する'\}\par

\noindent\textbf{カ：} \{'原典の主張の要約': '顧客が認めれば技術士補の名称で主体的に業務を行う', '独立判定': '×', '理由': '技術士補の名称で行えるのは技術士の補助としての業務である'\}\par

\noindent\textbf{キ：} \{'原典の主張の要約': '業務に関する知識技能を常に向上し専門領域を能力開発する', '独立判定': '○', '理由': '第47条の2の資質向上の責務に沿う'\}\par

\paragraph{問われる核心} 技術士法の義務・責務と適用条件

\paragraph{必要知識} 第44条の信用失墜禁止、第45条の守秘、第45条の2の公益、第46・47条の名称表示、第47条の2の継続的な資質向上。

\paragraph{考え方（査読待ち初稿）} アは○：顧客の業務情報を保護する原則として適切。第45条の正当な理由という例外まで否定するものとしては読まない。 イは×：守秘は退職・資格喪失後も存続する。 ウは×：公益と法令に反する指示まで受け入れない。 エは×：公衆の危険を顧客利益だけで放置しない。 オは×：第47条2項は名称表示の第46条を技術士補へ準用する。 カは×：技術士補の名称で行えるのは技術士の補助としての業務である。 キは○：第47条の2の資質向上の責務に沿う。 よって誤りは5件である。

\paragraph{正答を選ぶ理由} 誤りと判定する記述はイウエオカで、5件である。  よって誤りは5件である。

\Needspace{4\baselineskip}\paragraph{選択肢の検討（査読待ち初稿）}\mbox{}\par

\noindent\textbf{1：不適切な解答候補} 7件。選択肢は7件だが、誤りと判定したイウエオカの5件と合わない。\par

\noindent\textbf{2：不適切な解答候補} 6件。選択肢は6件だが、誤りと判定したイウエオカの5件と合わない。\par

\noindent\textbf{3：適切な解答候補} 5件。個別判定を数えると誤りは5件であり、この数に一致する。\par

\noindent\textbf{4：不適切な解答候補} 4件。選択肢は4件だが、誤りと判定したイウエオカの5件と合わない。\par

\noindent\textbf{5：不適切な解答候補} 3件。選択肢は3件だが、誤りと判定したイウエオカの5件と合わない。\par

\paragraph{再利用する考え方} 主語、名称表示をする場合という条件、退職後も残る義務を確認する。条文の名称と努力義務の性質を区別する。

\paragraph{共通理論} \hyperref[theory:engineer_act:duties]{第\ref*{theory:engineer_act:duties}節}

\par 参考：\cite{engineer_act_current}。

\par 難易度 2（校準初稿）、優先度未評価。

\subsection{令和4年度 問1}
\label{q:r04-reg-apt-01}

\noindent\textbf{原典：} sources/R04/aptitude.pdf、PDF物理1頁。

\textbf{公式正答：} 4（single）。

\paragraph{問題の要旨} アから順に×○×○×××となる。

\Needspace{4\baselineskip}\paragraph{確認した条件・記述（独自要約）}\mbox{}\par

\noindent\textbf{ア：} \{'原典の主張の要約': '退職後の守秘制約はない', '独立判定': '×', '理由': '守秘は資格喪失後も存続する'\}\par

\noindent\textbf{イ：} \{'原典の主張の要約': '名称表示している専門領域を能力開発し業務領域を拡大', '独立判定': '○', '理由': '継続的能力開発は専門領域の力量を高める。未登録部門の名称使用とは別である'\}\par

\noindent\textbf{ウ：} \{'原典の主張の要約': 'あらゆる顧客指示をそのまま実施', '独立判定': '×', '理由': '公益・法令に反する指示を無条件に受け入れない'\}\par

\noindent\textbf{エ：} \{'原典の主張の要約': '名称表示時は登録部門を明示し未登録部門を表示しない', '独立判定': '○', '理由': '第46条の条件に一致する'\}\par

\noindent\textbf{オ：} \{'原典の主張の要約': '顧客利益を害する場合は公益より守秘を優先', '独立判定': '×', '理由': '公衆の危険を顧客利益だけで放置しない'\}\par

\noindent\textbf{カ：} \{'原典の主張の要約': '顧客が認めれば技術士補の名称で主体的に業務を行う', '独立判定': '×', '理由': '第47条の補助という条件に反する'\}\par

\noindent\textbf{キ：} \{'原典の主張の要約': '登録部門以外の知識技能を重点的に向上', '独立判定': '×', '理由': '登録部門の知識技能も高める責務がある'\}\par

\paragraph{問われる核心} 技術士法の義務・責務と適用条件

\paragraph{必要知識} 第44条の信用失墜禁止、第45条の守秘、第45条の2の公益、第46・47条の名称表示、第47条の2の継続的な資質向上。

\paragraph{考え方（査読待ち初稿）} アは×：守秘は資格喪失後も存続する。 イは○：継続的能力開発は専門領域の力量を高める。未登録部門の名称使用とは別である。 ウは×：公益・法令に反する指示を無条件に受け入れない。 エは○：第46条の条件に一致する。 オは×：公衆の危険を顧客利益だけで放置しない。 カは×：第47条の補助という条件に反する。 キは×：登録部門の知識技能も高める責務がある。 判定列は×○×○×××である。

\paragraph{正答を選ぶ理由} アから順に×○×○×××となる。  判定列は×○×○×××である。

\Needspace{4\baselineskip}\paragraph{選択肢の検討（査読待ち初稿）}\mbox{}\par

\noindent\textbf{1：不適切な解答候補} アから順に×○××○×○。相違する箇所はエ（×とするが○）、オ（○とするが×）、キ（○とするが×）である。各記述の理由は独立判定に示した。\par

\noindent\textbf{2：不適切な解答候補} アから順に×××○×○×。相違する箇所はイ（×とするが○）、カ（○とするが×）である。各記述の理由は独立判定に示した。\par

\noindent\textbf{3：不適切な解答候補} アから順に○×○×○×○。相違する箇所はア（○とするが×）、イ（×とするが○）、ウ（○とするが×）、エ（×とするが○）、オ（○とするが×）、キ（○とするが×）である。各記述の理由は独立判定に示した。\par

\noindent\textbf{4：適切な解答候補} アから順に×○×○×××。全記述の判断に一致する。\par

\noindent\textbf{5：不適切な解答候補} アから順に○××○×○×。相違する箇所はア（○とするが×）、イ（×とするが○）、カ（○とするが×）である。各記述の理由は独立判定に示した。\par

\paragraph{再利用する考え方} 主語、名称表示をする場合という条件、退職後も残る義務を確認する。条文の名称と努力義務の性質を区別する。

\paragraph{共通理論} \hyperref[theory:engineer_act:duties]{第\ref*{theory:engineer_act:duties}節}

\par 参考：\cite{engineer_act_current}。

\par 難易度 2（校準初稿）、優先度未評価。

\subsection{令和6年度 問1}
\label{q:r06-reg-apt-01}

\noindent\textbf{原典：} sources/R06/aptitude.pdf、PDF物理1頁。

\textbf{公式正答：} 4（single）。

\paragraph{問題の要旨} 第4章の5空欄へ入る語を選ぶ。 正しい組合せは義務／公益／責務／技術部門／資質である。

\Needspace{4\baselineskip}\paragraph{確認した条件・記述（独自要約）}\mbox{}\par

\noindent\textbf{空欄順：} ['ア', 'イ', 'ウ', 'エ', 'オ']\par

\noindent\textbf{正しい語句の照合：} \{'ア': '義務', 'イ': '公益', 'ウ': '責務', 'エ': '技術部門', 'オ': '資質'\}\par

\noindent\textbf{出題時資料：} \par

\paragraph{問われる核心} 技術士法の義務・責務と適用条件

\paragraph{必要知識} 第44条の信用失墜禁止、第45条の守秘、第45条の2の公益、第46・47条の名称表示、第47条の2の継続的な資質向上。

\paragraph{考え方（査読待ち初稿）} 各空欄の主語と見出しから語句を特定する。 条文又は指定資料に一致する語句は、義務、公益、責務、技術部門、資質である。ほかの組合せは下記の相違がある。

\paragraph{正答を選ぶ理由} 第4章の5空欄へ入る語を選ぶ。 正しい組合せは義務／公益／責務／技術部門／資質である。 ほかの組合せは下記の相違がある。

\Needspace{4\baselineskip}\paragraph{選択肢の検討（査読待ち初稿）}\mbox{}\par

\noindent\textbf{1：不適切な解答候補} 責務／安全／義務／専門部門／能力。5語すべてが条文と異なる。\par

\noindent\textbf{2：不適切な解答候補} 義務／安全／責務／技術部門／資質。イは安全だけでなく公益である。\par

\noindent\textbf{3：不適切な解答候補} 責務／安全／義務／技術部門／能力。ア・イ・ウ・オが異なる。\par

\noindent\textbf{4：適切な解答候補} 義務／公益／責務／技術部門／資質。秘密保持の義務、公益・資質向上の責務、登録の技術部門に一致する。\par

\noindent\textbf{5：不適切な解答候補} 義務／公益／責務／専門部門／資質。エは専門部門ではなく技術部門である。\par

\paragraph{再利用する考え方} 主語、名称表示をする場合という条件、退職後も残る義務を確認する。条文の名称と努力義務の性質を区別する。

\paragraph{共通理論} \hyperref[theory:engineer_act:duties]{第\ref*{theory:engineer_act:duties}節}

\par 参考：\cite{engineer_act_current}。

\par 難易度 1（校準初稿）、優先度未評価。

\subsection{令和7年度 問1}
\label{q:r07-reg-apt-01}

\noindent\textbf{原典：} sources/R07/aptitude.pdf、PDF物理1頁。

\textbf{公式正答：} 3（single）。

\paragraph{問題の要旨} アから順に○○○×となる。

\Needspace{4\baselineskip}\paragraph{確認した条件・記述（独自要約）}\mbox{}\par

\noindent\textbf{ア：} \{'原典の主張の要約': '信用失墜や技術士等全体の不名誉となる行為を禁止', '独立判定': '○', '理由': '第44条の規定に一致する'\}\par

\noindent\textbf{イ：} \{'原典の主張の要約': '正当な理由なく秘密を漏らし又は盗用しない。資格を失った後も同じ', '独立判定': '○', '理由': '第45条の規定に一致する'\}\par

\noindent\textbf{ウ：} \{'原典の主張の要約': '公共安全・環境その他公益を害さない努力', '独立判定': '○', '理由': '第45条の2の規定に一致する'\}\par

\noindent\textbf{エ：} \{'原典の主張の要約': '未登録部門は他部門技術士の指導下で3年経験すれば表示できる', '独立判定': '×', '理由': '第46条は未登録部門の表示を禁止する。指導下の経験だけで表示できるという例外はない'\}\par

\paragraph{問われる核心} 技術士法の義務・責務と適用条件

\paragraph{必要知識} 第44条の信用失墜禁止、第45条の守秘、第45条の2の公益、第46・47条の名称表示、第47条の2の継続的な資質向上。

\paragraph{考え方（査読待ち初稿）} アは○：第44条の規定に一致する。 イは○：第45条の規定に一致する。 ウは○：第45条の2の規定に一致する。 エは×：第46条は未登録部門の表示を禁止する。指導下の経験だけで表示できるという例外はない。 判定列は○○○×である。

\paragraph{正答を選ぶ理由} アから順に○○○×となる。  判定列は○○○×である。

\Needspace{4\baselineskip}\paragraph{選択肢の検討（査読待ち初稿）}\mbox{}\par

\noindent\textbf{1：不適切な解答候補} アから順に○○××。相違する箇所はウ（×とするが○）である。各記述の理由は独立判定に示した。\par

\noindent\textbf{2：不適切な解答候補} アから順に××○○。相違する箇所はア（×とするが○）、イ（×とするが○）、エ（○とするが×）である。各記述の理由は独立判定に示した。\par

\noindent\textbf{3：適切な解答候補} アから順に○○○×。全記述の判断に一致する。\par

\noindent\textbf{4：不適切な解答候補} アから順に×××○。相違する箇所はア（×とするが○）、イ（×とするが○）、ウ（×とするが○）、エ（○とするが×）である。各記述の理由は独立判定に示した。\par

\noindent\textbf{5：不適切な解答候補} アから順に○×××。相違する箇所はイ（×とするが○）、ウ（×とするが○）である。各記述の理由は独立判定に示した。\par

\paragraph{再利用する考え方} 主語、名称表示をする場合という条件、退職後も残る義務を確認する。条文の名称と努力義務の性質を区別する。

\paragraph{共通理論} \hyperref[theory:engineer_act:duties]{第\ref*{theory:engineer_act:duties}節}

\par 参考：\cite{engineer_act_current}。

\par 難易度 2（校準初稿）、優先度未評価。

\section{公益・説明責任}

\subsection{令和5年度 問1}
\label{q:r05-reg-apt-01}

\noindent\textbf{原典：} sources/R05/aptitude.pdf、PDF物理1頁。

\textbf{公式正答：} 3（single）。

\paragraph{問題の要旨} アから順に×○××となる。

\Needspace{4\baselineskip}\paragraph{確認した条件・記述（独自要約）}\mbox{}\par

\noindent\textbf{ア：} \{'原典の主張の要約': '依頼者情報から独自調査した情報は守秘の対象外', '独立判定': '×', '理由': '独自調査だけで秘密性や契約上の権利は消えない'\}\par

\noindent\textbf{イ：} \{'原典の主張の要約': '自分や組織に不利な情報でも隠蔽せず公開に努める', '独立判定': '○', '理由': '公衆に必要な情報を組織の不利益だけで隠さない。ただし秘密を無条件に全公開する一般則ではない'\}\par

\noindent\textbf{ウ：} \{'原典の主張の要約': '公衆の安全に必要な情報を組織が公開しないときは守秘を優先', '独立判定': '×', '理由': '事態が改善しない場合にも公衆の安全を優先して適切な対応を検討する'\}\par

\noindent\textbf{エ：} \{'原典の主張の要約': '顧客判断に安全懸念があっても発言は不要', '独立判定': '×', '理由': '専門家として予測される危険を説明する'\}\par

\paragraph{問われる核心} 技術士法の義務・責務と適用条件

\paragraph{必要知識} 守秘の対象と存続、公益優先、客観的判断、情報の権利と利益相反。

\paragraph{考え方（査読待ち初稿）} アは×：独自調査だけで秘密性や契約上の権利は消えない。 イは○：公衆に必要な情報を組織の不利益だけで隠さない。ただし秘密を無条件に全公開する一般則ではない。 ウは×：事態が改善しない場合にも公衆の安全を優先して適切な対応を検討する。 エは×：専門家として予測される危険を説明する。 判定列は×○××である。

\paragraph{正答を選ぶ理由} アから順に×○××となる。  判定列は×○××である。

\Needspace{4\baselineskip}\paragraph{選択肢の検討（査読待ち初稿）}\mbox{}\par

\noindent\textbf{1：不適切な解答候補} アから順に○×○×。相違する箇所はア（○とするが×）、イ（×とするが○）、ウ（○とするが×）である。各記述の理由は独立判定に示した。\par

\noindent\textbf{2：不適切な解答候補} アから順に○○××。相違する箇所はア（○とするが×）である。各記述の理由は独立判定に示した。\par

\noindent\textbf{3：適切な解答候補} アから順に×○××。全記述の判断に一致する。\par

\noindent\textbf{4：不適切な解答候補} アから順に××○○。相違する箇所はイ（×とするが○）、ウ（○とするが×）、エ（○とするが×）である。各記述の理由は独立判定に示した。\par

\noindent\textbf{5：不適切な解答候補} アから順に××○×。相違する箇所はイ（×とするが○）、ウ（○とするが×）である。各記述の理由は独立判定に示した。\par

\paragraph{再利用する考え方} 公益確保と守秘を対立する二択にせず、情報の性質・危険・適切な相手への説明を検討する。独自調査で権利が自動的に移るとは限らない。

\paragraph{共通理論} \hyperref[theory:engineer_act:public]{第\ref*{theory:engineer_act:public}節}

\par 参考：\cite{engineer_act_current,cpd_v3}。

\par 難易度 2（校準初稿）、優先度未評価。

\section{CPD・継続研鑚}

\subsection{平成23年度 問4}
\label{q:h23-reg-apt-04}

\noindent\textbf{原典：} sources/H23/aptitude.pdf、PDF物理5頁。

\textbf{公式正答：} 1|5（multiple\_accepted）。

\par\textbf{検討未確定：} 原典の5肢はすべて確認済み。エの行政罰という表現と努力責務不履行への第36条2項の適用は追加の当時資料による検討が必要。1・5の候補を保持し、公式の複数正答を単一正答へ変更しない。

\paragraph{問題の要旨} ア・イ・ウは正しく、エの法的な評価に解釈上の問題が残る。

\Needspace{4\baselineskip}\paragraph{確認した条件・記述（独自要約）}\mbox{}\par

\noindent\textbf{ア：} CPDを技術者の基礎教育・試験・実務・資格取得と続く生涯システムの一部に位置付ける\par

\noindent\textbf{イ：} 各団体独自のCPD方法を相互認証する試みがある\par

\noindent\textbf{ウ：} 学協会の質量チェックが資格継続の条件となる場合がある\par

\noindent\textbf{エ：} 資質向上を怠る場合、行政罰に相当し登録取消し又は名称使用停止が可能とする\par

\paragraph{問われる核心} CPD制度と法的措置の区別

\paragraph{必要知識} CPDは能力向上を目的とする継続研鑚であり、通常業務との区別、実施記録と証拠、制度ごとの運用を確認する。

\paragraph{考え方（査読待ち初稿）} アは継続的研鑚を資格取得後にも必要とする生涯の専門家育成に位置付ける。イは団体間の実績活用、ウは制度ごとの継続条件である。エは第36条2項の登録取消し等という行政上の処分と、行政罰という言い方、努力責務を怠った場合の評価を区別しないため、独立検討では読み方に応じ1又は5が候補となる。単一の解釈に確定しない。

\paragraph{正答を選ぶ理由} ア・イ・ウは正しく、エの法的な評価に解釈上の問題が残る。 単一の解釈に確定しない。

\Needspace{4\baselineskip}\paragraph{選択肢の検討（査読待ち初稿）}\mbox{}\par

\noindent\textbf{1：適切な解答候補} ○○○○。ア・イ・ウは適切。エは登録取消し等を行政上の措置として広く読む場合の候補である。\par

\noindent\textbf{2：不適切な解答候補} ×○○○。アの生涯システムの位置付けを誤りとしており不適切である。\par

\noindent\textbf{3：不適切な解答候補} ×○×○。ア・ウを誤りとしており不適切である。\par

\noindent\textbf{4：不適切な解答候補} ○×○×。イの制度間協力を誤りとしており不適切である。\par

\noindent\textbf{5：適切な解答候補} ○○○×。ア・イ・ウは適切。エの行政罰という用語又は資質向上不履行と措置の結び付けを厳密に問題とする場合の候補である。\par

\paragraph{再利用する考え方} 勤務時間の経過と能力の向上は別である。学んだ内容と証拠を記録し、対象年度のガイドラインを参照する。

\paragraph{共通理論} \hyperref[theory:engineer_act:cpd]{第\ref*{theory:engineer_act:cpd}節}

\par 参考：\cite{engineer_act_current,cpd_v1}。

\par 難易度 2（校準初稿）、優先度未評価。

\subsection{平成24年度 問14}
\label{q:h24-reg-apt-14}

\noindent\textbf{原典：} sources/H24/aptitude.pdf、PDF物理12頁。

\textbf{公式正答：} 4（single）。

\paragraph{問題の要旨} アから順に○○×○となる。

\Needspace{4\baselineskip}\paragraph{確認した条件・記述（独自要約）}\mbox{}\par

\noindent\textbf{ア：} \{'原典の主張の要約': '業務に関する知識技能を常に向上させる', '独立判定': '○', '理由': '第47条の2の継続的資質向上に沿う'\}\par

\noindent\textbf{イ：} \{'原典の主張の要約': '学協会が質・量を確認し資格継続を制約する場合がある', '独立判定': '○', '理由': '制度ごとの運用であり全資格一律という主張ではない'\}\par

\noindent\textbf{ウ：} \{'原典の主張の要約': '通常の専門業務に勤務すればCPDの全要件を満たす', '独立判定': '×', '理由': '同じ知識を使う通常業務だけでは能力向上を示せない'\}\par

\noindent\textbf{エ：} \{'原典の主張の要約': '記録・証明がなければ実施していてもCPDとして認められない場合がある', '独立判定': '○', '理由': '実施内容と成果の証拠を残して説明できることが必要'\}\par

\paragraph{問われる核心} 技術士法の義務・責務と適用条件

\paragraph{必要知識} CPDは能力向上を目的とする継続研鑚であり、通常業務との区別、実施記録と証拠、制度ごとの運用を確認する。

\paragraph{考え方（査読待ち初稿）} アは○：第47条の2の継続的資質向上に沿う。 イは○：制度ごとの運用であり全資格一律という主張ではない。 ウは×：同じ知識を使う通常業務だけでは能力向上を示せない。 エは○：実施内容と成果の証拠を残して説明できることが必要。 判定列は○○×○である。

\paragraph{正答を選ぶ理由} アから順に○○×○となる。  判定列は○○×○である。

\Needspace{4\baselineskip}\paragraph{選択肢の検討（査読待ち初稿）}\mbox{}\par

\noindent\textbf{1：不適切な解答候補} アから順に○○○○。相違する箇所はウ（○とするが×）である。各記述の理由は独立判定に示した。\par

\noindent\textbf{2：不適切な解答候補} アから順に×○○×。相違する箇所はア（×とするが○）、ウ（○とするが×）、エ（×とするが○）である。各記述の理由は独立判定に示した。\par

\noindent\textbf{3：不適切な解答候補} アから順に×○×○。相違する箇所はア（×とするが○）である。各記述の理由は独立判定に示した。\par

\noindent\textbf{4：適切な解答候補} アから順に○○×○。全記述の判断に一致する。\par

\noindent\textbf{5：不適切な解答候補} アから順に○××○。相違する箇所はイ（×とするが○）である。各記述の理由は独立判定に示した。\par

\paragraph{再利用する考え方} 勤務時間の経過と能力の向上は別である。学んだ内容と証拠を記録し、対象年度のガイドラインを参照する。

\paragraph{共通理論} \hyperref[theory:engineer_act:cpd]{第\ref*{theory:engineer_act:cpd}節}

\par 参考：\cite{cpd_v1}。

\par 難易度 2（校準初稿）、優先度未評価。

\subsection{平成27年度 問2}
\label{q:h27-reg-apt-02}

\noindent\textbf{原典：} sources/H27/aptitude.pdf、PDF物理2頁。

\textbf{公式正答：} 5（single）。

\paragraph{問題の要旨} 平成26年4月のCPDガイドライン第2版の4空欄を選ぶ。 正しい組合せは自己研鑚／立場／記録／日頃従事である。

\Needspace{4\baselineskip}\paragraph{確認した条件・記述（独自要約）}\mbox{}\par

\noindent\textbf{空欄順：} ['ア', 'イ', 'ウ', 'エ']\par

\noindent\textbf{正しい語句の照合：} \{'ア': '自己研鑚', 'イ': '立場', 'ウ': '記録', 'エ': '日頃従事'\}\par

\noindent\textbf{出題時資料：} 出題指定の第2版2.2節で確認する。\par

\paragraph{問われる核心} 技術士法の義務・責務と適用条件

\paragraph{必要知識} CPDは能力向上を目的とする継続研鑚であり、通常業務との区別、実施記録と証拠、制度ごとの運用を確認する。

\paragraph{考え方（査読待ち初稿）} 各空欄の主語と見出しから語句を特定する。出題指定の第2版2.2節で確認する。 条文又は指定資料に一致する語句は、自己研鑚、立場、記録、日頃従事である。ほかの組合せは下記の相違がある。

\paragraph{正答を選ぶ理由} 平成26年4月のCPDガイドライン第2版の4空欄を選ぶ。 正しい組合せは自己研鑚／立場／記録／日頃従事である。 ほかの組合せは下記の相違がある。

\Needspace{4\baselineskip}\paragraph{選択肢の検討（査読待ち初稿）}\mbox{}\par

\noindent\textbf{1：不適切な解答候補} 自己研鑚／環境／記録／専門的に従事。イは立場、エは日頃従事である。\par

\noindent\textbf{2：不適切な解答候補} 自己学習／環境／掲載／日頃従事。アは自己研鑚、イは立場、ウは記録である。\par

\noindent\textbf{3：不適切な解答候補} 自己学習／立場／記録／専門的に従事。アは自己研鑚、エは日頃従事である。\par

\noindent\textbf{4：不適切な解答候補} 自己研鑚／環境／掲載／日頃従事。イは立場、ウは記録である。\par

\noindent\textbf{5：適切な解答候補} 自己研鑚／立場／記録／日頃従事。第2版2.2節の語句に一致する。\par

\paragraph{再利用する考え方} 勤務時間の経過と能力の向上は別である。学んだ内容と証拠を記録し、対象年度のガイドラインを参照する。

\paragraph{共通理論} \hyperref[theory:engineer_act:cpd]{第\ref*{theory:engineer_act:cpd}節}

\par 参考：\cite{cpd_v2}。

\par 難易度 1（校準初稿）、優先度未評価。

\subsection{平成28年度 問2}
\label{q:h28-reg-apt-02}

\noindent\textbf{原典：} sources/H28/aptitude.pdf、PDF物理2頁。

\textbf{公式正答：} 3（single）。

\paragraph{問題の要旨} アから順に○○×○となる。

\Needspace{4\baselineskip}\paragraph{確認した条件・記述（独自要約）}\mbox{}\par

\noindent\textbf{ア：} \{'原典の主張の要約': '記録・証明がなければCPDとして認められない場合がある', '独立判定': '○', '理由': '実施内容の記録と証拠による説明が必要'\}\par

\noindent\textbf{イ：} \{'原典の主張の要約': '業務に関する知識技能を常に向上させる', '独立判定': '○', '理由': '第47条の2の資質向上に沿う'\}\par

\noindent\textbf{ウ：} \{'原典の主張の要約': '通常の専門業務に勤務すればCPDの全要件を満たす', '独立判定': '×', '理由': '日常業務と、その業務に関連する新しい調査研究・能力向上を区別する'\}\par

\noindent\textbf{エ：} \{'原典の主張の要約': '学協会が質・量を確認し資格継続を制約する場合がある', '独立判定': '○', '理由': '制度ごとにチェックと継続条件を設ける場合がある'\}\par

\paragraph{問われる核心} 技術士法の義務・責務と適用条件

\paragraph{必要知識} CPDは能力向上を目的とする継続研鑚であり、通常業務との区別、実施記録と証拠、制度ごとの運用を確認する。

\paragraph{考え方（査読待ち初稿）} アは○：実施内容の記録と証拠による説明が必要。 イは○：第47条の2の資質向上に沿う。 ウは×：日常業務と、その業務に関連する新しい調査研究・能力向上を区別する。 エは○：制度ごとにチェックと継続条件を設ける場合がある。 判定列は○○×○である。

\paragraph{正答を選ぶ理由} アから順に○○×○となる。  判定列は○○×○である。

\Needspace{4\baselineskip}\paragraph{選択肢の検討（査読待ち初稿）}\mbox{}\par

\noindent\textbf{1：不適切な解答候補} アから順に×○×○。相違する箇所はア（×とするが○）である。各記述の理由は独立判定に示した。\par

\noindent\textbf{2：不適切な解答候補} アから順に×○○×。相違する箇所はア（×とするが○）、ウ（○とするが×）、エ（×とするが○）である。各記述の理由は独立判定に示した。\par

\noindent\textbf{3：適切な解答候補} アから順に○○×○。全記述の判断に一致する。\par

\noindent\textbf{4：不適切な解答候補} アから順に○○××。相違する箇所はエ（×とするが○）である。各記述の理由は独立判定に示した。\par

\noindent\textbf{5：不適切な解答候補} アから順に○×○○。相違する箇所はイ（×とするが○）、ウ（○とするが×）である。各記述の理由は独立判定に示した。\par

\paragraph{再利用する考え方} 勤務時間の経過と能力の向上は別である。学んだ内容と証拠を記録し、対象年度のガイドラインを参照する。

\paragraph{共通理論} \hyperref[theory:engineer_act:cpd]{第\ref*{theory:engineer_act:cpd}節}

\par 参考：\cite{cpd_v2}。

\par 難易度 2（校準初稿）、優先度未評価。

\subsection{平成30年度 問3}
\label{q:h30-reg-apt-03}

\noindent\textbf{原典：} sources/H30/aptitude.pdf、PDF物理3頁。

\textbf{公式正答：} 3（single）。

\paragraph{問題の要旨} 平成29年4月のCPDガイドライン第3版の4空欄を選ぶ。 正しい組合せは継続研鑚／専門家／立場／記録である。

\Needspace{4\baselineskip}\paragraph{確認した条件・記述（独自要約）}\mbox{}\par

\noindent\textbf{空欄順：} ['ア', 'イ', 'ウ', 'エ']\par

\noindent\textbf{正しい語句の照合：} \{'ア': '継続研鑚', 'イ': '専門家', 'ウ': '立場', 'エ': '記録'\}\par

\noindent\textbf{出題時資料：} 出題指定の第3版2.2節で確認する。\par

\paragraph{問われる核心} 技術士法の義務・責務と適用条件

\paragraph{必要知識} CPDは能力向上を目的とする継続研鑚であり、通常業務との区別、実施記録と証拠、制度ごとの運用を確認する。

\paragraph{考え方（査読待ち初稿）} 各空欄の主語と見出しから語句を特定する。出題指定の第3版2.2節で確認する。 条文又は指定資料に一致する語句は、継続研鑚、専門家、立場、記録である。ほかの組合せは下記の相違がある。

\paragraph{正答を選ぶ理由} 平成29年4月のCPDガイドライン第3版の4空欄を選ぶ。 正しい組合せは継続研鑚／専門家／立場／記録である。 ほかの組合せは下記の相違がある。

\Needspace{4\baselineskip}\paragraph{選択肢の検討（査読待ち初稿）}\mbox{}\par

\noindent\textbf{1：不適切な解答候補} 継続学習／技術者／環境／記録。ア・イ・ウが出題指定資料と異なる。\par

\noindent\textbf{2：不適切な解答候補} 継続学習／専門家／環境／記載。ア・ウ・エが異なる。\par

\noindent\textbf{3：適切な解答候補} 継続研鑚／専門家／立場／記録。第3版2.2節の語句に一致する。\par

\noindent\textbf{4：不適切な解答候補} 継続学習／技術者／環境／記載。ア・イ・ウ・エが異なる。\par

\noindent\textbf{5：不適切な解答候補} 継続研鑚／専門家／立場／記載。エは記載ではなく記録である。\par

\paragraph{再利用する考え方} 勤務時間の経過と能力の向上は別である。学んだ内容と証拠を記録し、対象年度のガイドラインを参照する。

\paragraph{共通理論} \hyperref[theory:engineer_act:cpd]{第\ref*{theory:engineer_act:cpd}節}

\par 参考：\cite{cpd_v3}。

\par 難易度 1（校準初稿）、優先度未評価。

\subsection{令和6年度 問15}
\label{q:r06-reg-apt-15}

\noindent\textbf{原典：} sources/R06/aptitude.pdf、PDF物理15頁。

\textbf{公式正答：} 2（single）。

\paragraph{問題の要旨} 正しいと判定する記述はアイウオで、4件である。

\Needspace{4\baselineskip}\paragraph{確認した条件・記述（独自要約）}\mbox{}\par

\noindent\textbf{ア：} \{'原典の主張の要約': 'CPDで業務知識技能を常に向上', '独立判定': '○', '理由': '資質向上の責務に沿う'\}\par

\noindent\textbf{イ：} \{'原典の主張の要約': '実施内容を記録・証拠とともに第三者へ説明可能にする', '独立判定': '○', '理由': '本人による実施記録と証明が必要'\}\par

\noindent\textbf{ウ：} \{'原典の主張の要約': '2021年省令改正で登録簿に取組状況欄を追加しCPD実績記載・認定・公表制度を実施', '独立判定': '○', '理由': '2021年9月の省令改正と新制度で確認できる。全技術士に一律の資格更新が課されたという意味ではない'\}\par

\noindent\textbf{エ：} \{'原典の主張の要約': '通常の既知の専門業務に勤務すればCPDの全要件を満たす', '独立判定': '×', '理由': '通常業務の遂行と知識・力量の向上を区別する'\}\par

\noindent\textbf{オ：} \{'原典の主張の要約': '学協会が質・量を確認し資格継続を制約する場合がある', '独立判定': '○', '理由': '学協会・資格制度ごとの継続条件がある'\}\par

\paragraph{問われる核心} 技術士法の義務・責務と適用条件

\paragraph{必要知識} CPDは能力向上を目的とする継続研鑚であり、通常業務との区別、実施記録と証拠、制度ごとの運用を確認する。

\paragraph{考え方（査読待ち初稿）} アは○：資質向上の責務に沿う。 イは○：本人による実施記録と証明が必要。 ウは○：2021年9月の省令改正と新制度で確認できる。全技術士に一律の資格更新が課されたという意味ではない。 エは×：通常業務の遂行と知識・力量の向上を区別する。 オは○：学協会・資格制度ごとの継続条件がある。 よって正しいは4件である。

\paragraph{正答を選ぶ理由} 正しいと判定する記述はアイウオで、4件である。  よって正しいは4件である。

\Needspace{4\baselineskip}\paragraph{選択肢の検討（査読待ち初稿）}\mbox{}\par

\noindent\textbf{1：不適切な解答候補} 5件。選択肢は5件だが、正しいと判定したアイウオの4件と合わない。\par

\noindent\textbf{2：適切な解答候補} 4件。個別判定を数えると正しいは4件であり、この数に一致する。\par

\noindent\textbf{3：不適切な解答候補} 3件。選択肢は3件だが、正しいと判定したアイウオの4件と合わない。\par

\noindent\textbf{4：不適切な解答候補} 2件。選択肢は2件だが、正しいと判定したアイウオの4件と合わない。\par

\noindent\textbf{5：不適切な解答候補} 1件。選択肢は1件だが、正しいと判定したアイウオの4件と合わない。\par

\paragraph{再利用する考え方} 勤務時間の経過と能力の向上は別である。学んだ内容と証拠を記録し、対象年度のガイドラインを参照する。

\paragraph{共通理論} \hyperref[theory:engineer_act:cpd]{第\ref*{theory:engineer_act:cpd}節}

\par 参考：\cite{cpd_new2021,cpd_v3}。

\par 難易度 2（校準初稿）、優先度未評価。
